% !TEX program = xelatex
% Example: Overlays with \pause and \onslide
% First frame reveals bullets one at a time with \pause.
% Second frame uses \onslide<2-> to control visibility by overlay number.
\documentclass{beamer}
\usetheme{default}

\title{Overlays and Animations}
\author{Dr.~Richard Feynman}
\institute{Caltech}
\date{2025}

\begin{document}

\begin{frame}
  \titlepage
\end{frame}

\begin{frame}{Revealing with \texttt{\textbackslash pause}}
  Each bullet appears on a new overlay:
  \begin{itemize}
    \item First, we state the problem.
    \pause
    \item Then, we introduce the method.
    \pause
    \item Next, we present the result.
    \pause
    \item Finally, we discuss the implications.
  \end{itemize}
\end{frame}

\begin{frame}{Revealing with \texttt{\textbackslash onslide}}
  \onslide<1->{This text is visible from slide 1 onward.}
  \vspace{0.5em}

  \onslide<2->{This text appears starting at overlay 2.}
  \vspace{0.5em}

  \onslide<3->{This text appears starting at overlay 3.}
  \vspace{0.5em}

  \onslide<2->{Note: the second line was hidden, then revealed.}
\end{frame}

\end{document}
